\documentclass[11pt,letterpaper]{article}
% Standalone source: the bibliography is embedded below.
% Compile with pdfLaTeX (two passes) or an equivalent LaTeX engine.
\usepackage[margin=1in]{geometry}
\usepackage[T1]{fontenc}
\usepackage[utf8]{inputenc}
\usepackage{lmodern}
\usepackage{microtype}
\usepackage{amsmath,amssymb,amsthm,mathtools}
\usepackage{booktabs,tabularx,array}
\usepackage{enumitem}
\usepackage{xcolor}
\usepackage{float}
\floatstyle{ruled}
\newfloat{algorithm}{tbp}{loa}
\floatname{algorithm}{Algorithm}
\usepackage[font=small,labelfont=bf,labelsep=period]{caption}
\usepackage[numbers,sort&compress]{natbib}
\usepackage{aliascnt}
\usepackage{hyperref}
\usepackage[nameinlink,capitalize,noabbrev]{cleveref}
\usepackage{url}

\definecolor{linkblue}{RGB}{25,65,105}
\hypersetup{
  colorlinks=true,linkcolor=linkblue,citecolor=linkblue,urlcolor=linkblue,
  pdftitle={Optimal Exponents for Short-Lived High-Volume Bandits},
  pdfauthor={Colombo Team},
  pdfsubject={Finite-feedback Bayesian bandits and matching loss exponents},
  pdfkeywords={Bayesian bandits, short-lived arms, batched feedback, minimax rates}
}
\numberwithin{equation}{section}
\newtheorem{theorem}{Theorem}[section]
\newaliascnt{proposition}{theorem}
\newtheorem{proposition}[proposition]{Proposition}
\aliascntresetthe{proposition}
\newaliascnt{lemma}{theorem}
\newtheorem{lemma}[lemma]{Lemma}
\aliascntresetthe{lemma}
\newaliascnt{corollary}{theorem}
\newtheorem{corollary}[corollary]{Corollary}
\aliascntresetthe{corollary}
\theoremstyle{definition}
\newaliascnt{definition}{theorem}
\newtheorem{definition}[definition]{Definition}
\aliascntresetthe{definition}
\newaliascnt{assumption}{theorem}
\newtheorem{assumption}[assumption]{Assumption}
\aliascntresetthe{assumption}
\theoremstyle{remark}
\newaliascnt{remark}{theorem}
\newtheorem{remark}[remark]{Remark}
\aliascntresetthe{remark}
\crefname{proposition}{Proposition}{Propositions}
\crefname{lemma}{Lemma}{Lemmas}
\crefname{corollary}{Corollary}{Corollaries}
\crefname{definition}{Definition}{Definitions}
\crefname{assumption}{Assumption}{Assumptions}
\crefname{remark}{Remark}{Remarks}
\Crefname{proposition}{Proposition}{Propositions}
\Crefname{lemma}{Lemma}{Lemmas}
\Crefname{corollary}{Corollary}{Corollaries}
\Crefname{definition}{Definition}{Definitions}
\Crefname{assumption}{Assumption}{Assumptions}
\Crefname{remark}{Remark}{Remarks}

\newcommand{\E}{\mathbb{E}}
\newcommand{\Pp}{\mathbb{P}}
\newcommand{\1}{\mathbf{1}}
\newcommand{\F}{\mathcal{F}}
\newcommand{\Hh}{\mathcal{H}}
\newcommand{\A}{\mathcal{A}}
\newcommand{\D}{\mathcal{D}}
\newcommand{\Loss}{\operatorname{Loss}}
\newcommand{\wtO}{\widetilde{O}}
\newcommand{\wtOmega}{\widetilde{\Omega}}
\newcommand{\wtTheta}{\widetilde{\Theta}}
\newcommand{\eps}{\varepsilon}
\newcommand{\BetaDist}{\operatorname{Beta}}
\newcommand{\Unif}{\operatorname{Unif}}
\newcommand{\Bin}{\operatorname{Bin}}
\newcommand{\Reg}{\operatorname{Reg}}
\newcolumntype{Y}{>{\raggedright\arraybackslash}X}
\setlength{\emergencystretch}{1.5em}
\setlength{\parskip}{2pt}
\setlist[enumerate]{itemsep=3pt,topsep=5pt}

\title{Optimal Exponents for\\Short-Lived High-Volume Bandits}
\author{Colombo Team\\[3pt]
  \small\href{mailto:colombo.2026@outlook.com}{\texttt{colombo.2026@outlook.com}}}
\date{}

\begin{document}
\maketitle

\begin{abstract}
We determine the optimal loss exponent in the Bernoulli short-lived
high-volume bandit model. In every round, $k$ arms arrive, each remains
available through age $w$, and a learner allocates $n$ plays before observing
any reward from that round. Arm means are independent draws from a prior
whose density is bounded above and below on $[0,1]$. For fixed $w\ge1$, define
$\alpha_w=(2^w-1)/(2^{w+1}-1)$. We construct a policy with long-run
Bayesian loss at most
$C_{w,c_1,c_2}\{\log(n)/k+(\log(n)/n)^{\alpha_w}\}$.
For the uniform prior, every policy incurs loss at least
$c_w\max\{k^{-1},n^{-\alpha_w}\}$.
Consequently, when $k=\lfloor n^\rho\rfloor$, the minimax loss over the
bounded-density class is
$\wtTheta(n^{-\min\{\rho,\alpha_w\}})$.
The policy screens each cohort at increasing sampling depths, retaining an
arm only when all its screening observations are successes; deterministic
caps enforce the budget on every sample path. The lower bound rests on a
posterior-discovery inequality: after $j$ feedback batches, the expected
number of arms whose posterior shortfall ever falls below $x$ is at most
$C_j Bx^{2^j/(2^j-1)}$ for sufficiently small $x$, where $B$ is the
expected Bayesian endpoint cost.
This inequality accommodates adaptive allocations and overlapping cohorts.
The result sharpens the previously reported upper exponent
$w/[2(w+1)]$ and lower exponent $1/2$ to a matching exponent.
The remaining gap to $1/2$ is exactly
$1/[2(2^{w+1}-1)]$, which decreases exponentially with the number of
available feedback opportunities.
\end{abstract}

\noindent\textbf{Keywords.}
Bayesian bandits; short-lived arms; batched feedback; exploration and
exploitation; minimax loss.

\begin{samepage}
\noindent\textbf{AI Usage.}
Nearly the entire research pipeline for this paper was carried out by
Colombo, powered by GPT-5.6 Sol.
\par
\end{samepage}

\section{Introduction}

Learning is difficult when promising options disappear before they can be
evaluated. In online experimentation and content recommendation, a large
number of items may arrive between successive model updates, while each item
remains useful for only a few updates. A large sampling budget does not
remove this difficulty: many observations can be collected at once, but the
learner has only a few opportunities to use them before an item expires.
The short-lived high-volume bandit (SLHVB) model of
\citet{jia2023short,jia2026short} makes this interaction between sampling
volume and feedback explicit.

In the Bernoulli version considered here, $k$ new arms arrive in every round,
and their means are independent draws from a prior on $[0,1]$. Each round
contains exactly $n$ plays, with repeated plays of the same arm allowed.
The entire allocation is fixed before that round's rewards are revealed.
Following the formal indexing of \citet{jia2023short}, an arm born in round
$t$ remains available in rounds $t,\ldots,t+w$. It therefore has at most
$w$ completed feedback batches before its final use. The performance
criterion is average loss per play relative to the best arm currently alive.

For fixed $w$ and $k=\lfloor n^\rho\rfloor$, the previously reported
upper and lower bounds are, respectively,
\begin{equation}\label{eq:previous-rates}
  \wtO\!\left(n^{-\min\{\rho,\theta_w\}}\right)
  \quad\text{and}\quad
  \Omega\!\left(n^{-\min\{\rho,1/2\}}\right),
  \qquad \theta_w:=\frac{w}{2(w+1)}.
\end{equation}
These bounds coincide up to logarithmic factors when
$\rho\le\theta_w$, but leave a polynomial gap beyond that regime
\citep[Theorems 3.3 and 4.10]{jia2023short}.
The journal version reports the same exponents \citep{jia2026short}.
The unresolved question is how the finite number of feedback opportunities
changes the optimal rate when many fresh arms are available.

\subsection{Main result and contributions}

We identify the matching exponent
\begin{equation}\label{eq:alpha-beta-def}
  \alpha_w:=\frac{2^w-1}{2^{w+1}-1},
  \qquad
  \beta_w:=\frac{2^w}{2^w-1},
  \qquad
  \alpha_w=\frac{1}{1+\beta_w}.
\end{equation}
For fixed density bounds $0<c_1\le f\le c_2$ and fixed $w\ge1$,
we prove
\begin{equation}\label{eq:intro-main}
  c_w\max\{k^{-1},n^{-\alpha_w}\}
  \ \le\ \Loss_n^*
  \ \le\
  C_{w,c_1,c_2}\left\{
    \frac{\log n}{k}+
    \left(\frac{\log n}{n}\right)^{\alpha_w}
  \right\}.
\end{equation}
Here $\Loss_n^*$ denotes the minimax Bayesian loss over the bounded-density
prior class; the precise definition and quantifiers appear in
\cref{sec:model}. The lower bound already holds for the uniform prior, and
the upper bound uses a single policy for the whole class.

The two terms have distinct meanings. The term $k^{-1}$ reflects the arrival
of unobserved arms: the best fresh arm may outperform every previously
observed arm. The term $n^{-\alpha_w}$ measures the cost of finding a good
arm with only $w$ feedback opportunities. Their balance occurs at
$k\asymp n^{\alpha_w}$ at the level of polynomial exponents.
For $w=1,2,3,4$, the exponents are
$1/3,3/7,7/15,15/31$, respectively. In particular,
\begin{equation}\label{eq:exponent-gap}
  \frac12-\alpha_w=\frac{1}{2(2^{w+1}-1)}.
\end{equation}
This is an exact identity for the fixed-lifetime exponents; our constants
are not uniform when $w$ grows with $n$.

The proof has three components.
\begin{enumerate}[leftmargin=2em]
  \item \textbf{A capped all-success sieve.}
  Each cohort is screened at progressively larger depths. The depths
  equalize the expected loss spent at each age, while deterministic caps
  guarantee that every round uses at most its exploration budget.
  The policy needs the density bounds but not the density itself.
  \item \textbf{A finite-batch discovery inequality.}
  We bound the expected number of arms whose posterior endpoint shortfall
  reaches a prescribed threshold in $j$ batches in terms of the Bayesian
  cost spent on those arms. The exponent is
  $\beta_j=2^j/(2^j-1)$. This bound permits arbitrary nonanticipating
  allocations and observations from independent outside arms.
  \item \textbf{A streaming lower bound in every arrival regime.}
  We embed the alive cohorts into a $w$-batch experiment and charge each
  past cost to at most $w$ future windows. A separate, explicit
  $\Omega_w(k^{-1})$ lower bound converts endpoint shortfall to the actual
  loss against the best alive arm, including the case $w=1$.
\end{enumerate}

\begin{table}[t]
\centering
\caption{Quantitative comparison in the Bernoulli SLHVB model. Throughout,
$w\ge1$ and $\rho>0$ are fixed and $k=\lfloor n^\rho\rfloor$.
Upper bounds are achievable losses; lower bounds apply to every policy.
The earlier entries are the guarantees reported by
\citet{jia2023short,jia2026short}.}
\label{tab:rates}
\small
\renewcommand{\arraystretch}{1.28}
\begin{tabularx}{\linewidth}{@{}p{0.23\linewidth}p{0.18\linewidth}Y@{}}
\toprule
Result & Prior condition & Loss guarantee \\
\midrule
Previous upper bound & Bounded density
  & $\wtO(n^{-\min\{\rho,w/[2(w+1)]\}})$ \\
Previous lower bound & Bounded density
  & $\Omega(n^{-\min\{\rho,1/2\}})$ \\
Present upper bound & Bounded density
  & $O_{w,c_1,c_2}\!\left(\log(n)/k+
       (\log(n)/n)^{\alpha_w}\right)$ \\
Present lower bound & Uniform prior
  & $\Omega_w(\max\{k^{-1},n^{-\alpha_w}\})$ \\
\bottomrule
\end{tabularx}
\end{table}

\begin{table}[t]
\centering
\caption{The exponent of loss in each arrival regime, with
$\theta_w=w/[2(w+1)]<\alpha_w<1/2$. An upper exponent $\gamma$
denotes $\wtO(n^{-\gamma})$; a lower exponent denotes
$\Omega(n^{-\gamma})$. The last column is the matching minimax exponent
over the bounded-density class.}
\label{tab:regimes}
\small
\renewcommand{\arraystretch}{1.28}
\begin{tabular}{@{}lccc@{}}
\toprule
Arrival regime & Previous upper & Previous lower & Present exponent \\
\midrule
$0<\rho\le\theta_w$       & $\rho$     & $\rho$ & $\rho$ \\
$\theta_w<\rho\le\alpha_w$ & $\theta_w$ & $\rho$ & $\rho$ \\
$\alpha_w<\rho<1/2$       & $\theta_w$ & $\rho$ & $\alpha_w$ \\
$\rho\ge1/2$             & $\theta_w$ & $1/2$  & $\alpha_w$ \\
\bottomrule
\end{tabular}
\end{table}

\subsection{Related work}

\paragraph{Streaming arms with finite lifetimes.}
\citet{jia2023short} introduce the random-input SLHVB model and analyze a
layered sieve based on batched successive elimination. Its reported
guarantees give the interval in \cref{eq:previous-rates}; the journal
article develops the model and its applications further
\citep{jia2026short}. Our result sharpens this interval for Bernoulli
observations. The improvement uses their vanishing variance near mean one
and does not establish the same rate for arbitrary sub-Gaussian reward
families. The lower bound replaces an exploration and cold-start argument
with a quantitative restriction on posterior discovery under a fixed
number of updates.

\paragraph{Batched bandits and limited adaptivity.}
\citet{perchet2016batched} establish, up to logarithmic factors, the
two-arm minimax cumulative-regret rate
$T^{1/(2-2^{1-M})}$ for $M$ batches.
\citet{gao2019batched} extend the analysis to multiple arms and
adaptive batch sizes. After division by $T$ and substitution of
$M=w+1$, the two-arm exponent satisfies
\begin{equation}\label{eq:batch-identity}
  1-\frac{1}{2-2^{-w}}=\alpha_w.
\end{equation}
This numerical correspondence does not provide a reduction: those
problems have a fixed arm set and worst-case means, whereas SLHVB has
random arrivals, expiration, and a changing comparator.
The posterior-discovery argument handles these additional features.
Limited-round identification problems studied by
\citet{agarwal2017limited} give another perspective on the cost of
restricted feedback, with a different objective.

\paragraph{Infinite reservoirs and many arms.}
The distribution of means near the best possible value is central to
infinite-arm bandits \citep{berry1997infinite,wang2008infinite}.
Success-run screening has a direct precedent in the Bernoulli policies
of \citet{bonald2013two}. Thus, the use of successes to identify good arms
is an established idea. Here the contribution is its scheduling across
a fixed lifetime, capacity control for simultaneous cohorts, and a
matching lower bound. Related work studies greedy policies in Bayesian
many-arm problems \citep{bayati2020greedy} and adaptive recommendation
from broad reservoir classes \citep{pilliat2026adaptive}.
In these reservoir models, an identified arm can remain available;
our policy must repeatedly replace it before expiration.

\paragraph{Mortality and reward decay.}
Mortal bandits \citep{chakrabarti2008mortal}, irrevocable bandits
\citep{farias2011irrevocable}, and rotting bandits
\citep{levine2017rotting} impose other restrictions on how arms may
be used over time. They capture disappearance, irreversible abandonment,
or deteriorating rewards, respectively. The present analysis concerns
exogenous cohort arrivals and deterministic calendar-time expiration,
with repeated plays and delayed feedback within each round.

\section{Model and main theorem}\label{sec:model}

\subsection{Arrivals, rewards, and admissible policies}

Fix integers $n,k\ge1$ and $w\ge1$. A cohort $A_t$ of $k$ distinct arms
arrives at the start of round $t\ge1$. Cohorts are disjoint, and an arm
born at $t$ is available through round $t+w$. The alive set is
\[
  \A_t:=\bigcup_{s=(t-w)\vee1}^{t}A_s.
\]
An arm's age is the number of rounds since its arrival. In particular,
there are $w+1$ alive cohorts after initialization and $w$ useful
feedback opportunities before an arm's final age-$w$ allocation.
This is the formal lifetime convention of \citet{jia2023short}.

Each mean $\mu_a$ is independently drawn from a distribution $D$ on
$[0,1]$. Conditional on all means, rewards are independent across
arms and plays, with reward law $\operatorname{Bernoulli}(\mu_a)$.
Write
\[
  X_a:=1-\mu_a
\]
for the arm's \emph{endpoint shortfall}. The prior may be known to
the learner; our upper-bound policy only uses bounds on its density.

Let $\F_t$ be the sigma-field generated by the arm identities, allocations, and rewards
revealed through round $t$, together with the policy's independent
random seed. Arrival times and identities carry no information about
the means. An admissible policy chooses nonnegative integer
multiplicities $\pi_t(a)$, measurable with respect to $\F_{t-1}$, with
\begin{equation}\label{eq:allocation}
  \sum_{a\in\A_t}\pi_t(a)=n.
\end{equation}
All rewards from round $t$ are revealed after this entire allocation
has been fixed. In particular, there is no stopping or adaptation
within a round.

\subsection{Loss and prior class}

Let $\mu_t^*:=\max_{a\in\A_t}\mu_a$. The expected average loss over
$T$ rounds is
\begin{equation}\label{eq:finite-loss}
  \Loss_{n,T}(\pi;D)
  :=\frac{1}{nT}\E_D\!\left[
    \sum_{t=1}^{T}\sum_{a\in\A_t}
      \pi_t(a)(\mu_t^*-\mu_a)
  \right].
\end{equation}
The expectation includes the means, rewards, and policy randomization.
We write
\[
  \overline{\Loss}_n(\pi;D)
  :=\limsup_{T\to\infty}\Loss_{n,T}(\pi;D),
  \qquad
  \underline{\Loss}_n(\pi;D)
  :=\liminf_{T\to\infty}\Loss_{n,T}(\pi;D).
\]
This extends the expected-average definition of \citet{jia2023short}
by using upper and lower limits, so a long-run limit need not exist
for an arbitrary policy. An alternative expectation-of-limsup
criterion is treated separately in \cref{app:long-run}.

\begin{assumption}[Bounded prior density]\label{ass:bounded-density}
The distribution $D$ has a density $f$ on $[0,1]$ satisfying
\[
  c_1\le f(u)\le c_2
  \quad\text{for almost every }u\in[0,1],
  \qquad 0<c_1\le1\le c_2<\infty.
\]
\end{assumption}

Denote this class by $\mathfrak D(c_1,c_2)$. It contains the uniform
distribution. The restriction $c_1\le1\le c_2$ is necessary for
nonemptiness. Throughout, $w,c_1,c_2$ are fixed. Constants in upper
bounds may depend on these three parameters; constants in uniform-prior
lower bounds may depend on $w$. The notation $\wtO$ and $\wtTheta$
suppresses powers of $\log n$, not powers of $n$.

\begin{theorem}[Matching finite-lifetime exponent]\label{thm:main}
Fix $w\ge1$ and $0<c_1\le1\le c_2<\infty$, and let $\alpha_w$
be as in \cref{eq:alpha-beta-def}. There exist
$C=C(w,c_1,c_2)<\infty$, $c=c(w)>0$, and
$n_0=n_0(w,c_1,c_2)$ such that, for all $n\ge n_0$ and all
integers $k\ge1$, the following statements hold.
\begin{enumerate}[label=(\roman*),leftmargin=2.2em]
  \item There is a single admissible policy $\pi$, depending on
  $(n,k,w,c_1,c_2)$ but not on the particular $D$, such that every
  $D\in\mathfrak D(c_1,c_2)$ satisfies
  \begin{equation}\label{eq:main-upper}
    \overline{\Loss}_n(\pi;D)
    \le C\left[
      \frac{\log n}{k}
      +\left(\frac{\log n}{n}\right)^{\alpha_w}
    \right].
  \end{equation}
  \item For $D=\Unif[0,1]$, every admissible policy, including one
  that knows this prior, satisfies
  \begin{equation}\label{eq:main-lower}
    \underline{\Loss}_n(\pi;D)
    \ge c\max\left\{\frac1k,n^{-\alpha_w}\right\}.
  \end{equation}
\end{enumerate}
When the right-hand side of \cref{eq:main-upper} is at least one,
the bound follows from $\overline{\Loss}_n\le1$.
\end{theorem}

Define
\[
  \Loss_n^*(c_1,c_2)
  :=\inf_\pi\sup_{D\in\mathfrak D(c_1,c_2)}
          \overline{\Loss}_n(\pi;D),
\]
where a policy is allowed to take $D$ as input. Allowing this knowledge
does not change either side of the theorem: the upper bound is
independent of $D$, and the lower bound holds for a fixed known prior.

\begin{corollary}[Polynomial arrivals]\label{cor:rho}
For fixed $\rho>0$ and $k=\lfloor n^\rho\rfloor$,
\[
  \Loss_n^*(c_1,c_2)
  =\wtTheta_{w,c_1,c_2}\!
    \left(n^{-\min\{\rho,\alpha_w\}}\right).
\]
The same exponent holds for
$\inf_\pi\overline{\Loss}_n(\pi;\Unif[0,1])$.
\end{corollary}

\begin{proof}
The uniform prior belongs to the class, and
$\overline{\Loss}_n\ge\underline{\Loss}_n$. Thus
\cref{eq:main-lower} gives the lower bound in both formulations.
The single policy in \cref{eq:main-upper} gives the upper bound.
\end{proof}

\subsection{Proof overview}

The useful statistical event is a run of screening successes.
Under \cref{ass:bounded-density}, the probability of $s$ successes is
of order $s^{-1}$, and their posterior endpoint shortfall is also
of order $s^{-1}$. We start with
$m\asymp(\log n)/\eps$ candidates per cohort and choose cumulative depths
\[
  S_j\asymp\eps^{-a_j},
  \qquad
  a_j=\frac{2^j-1}{2^w-1},
  \qquad j=1,\ldots,w.
\]
The expected endpoint cost of level $j$ is of order
$mS_j/(S_{j-1}+1)^2$. The identity
$1+a_j-2a_{j-1}=\beta_w$ makes these costs equal in exponent.
After division by the round budget, exploration costs
$O((\log n)\eps^{-\beta_w}/n)$, while exploitation costs
$O(\eps)$. Balancing the two gives
$\eps\asymp(\log n/n)^{\alpha_w}$. Enough candidates are available
when $\eps\gtrsim\log(n)/k$.

For the lower bound, let $g$ be an arm's posterior mean shortfall.
A beta-posterior small-ball estimate bounds the chance that a future
posterior reaches $x$ by $O(x/g)$. Reaching that threshold also requires
at least order $x^{-1}$ observations. Combining the probability of
crossing with the cost of sampling deeply yields the recursion
\[
  \beta_1=2,
  \qquad
  \beta_j=\frac{2\beta_{j-1}}{\beta_{j-1}+1}.
\]
Its solution is the exponent in \cref{eq:alpha-beta-def}.
The remaining steps account for overlapping cohorts and subtract the
shortfall of the best alive arm without losing the arrival-limited regime.

\section{A feasible all-success screening policy}\label{sec:upper}

The upper-bound policy exploits a simple endpoint identity.  If an arm has
produced $s$ designated screening successes, then the probability of that
event and its posterior endpoint shortfall are both of order $1/(s+1)$.

\begin{lemma}[Success moments]\label{lem:success-moments}
Under \cref{ass:bounded-density}, for every integer $s\ge0$,
\begin{align}
  \frac{c_1}{s+1}
  &\le p_s:=\E[\mu^s]
  \le \frac{c_2}{s+1},                                    \label{eq:ps}\\
  q_s:=\E[(1-\mu)\mu^s]
  &\le \frac{c_2}{(s+1)(s+2)}.                            \label{eq:qs}
\end{align}
Consequently,
\begin{equation}\label{eq:posterior-success-mean}
  \E[1-\mu\mid s\text{ screening successes}]
  =\frac{q_s}{p_s}
  \le \frac{c_2/c_1}{s+2}.
\end{equation}
\end{lemma}

\begin{proof}
Integrating against $f$ and using $c_1\le f\le c_2$ gives
\[
  c_1\int_0^1u^s\,du
  \le \E[\mu^s]
  \le c_2\int_0^1u^s\,du,
\]
and
\[
  \E[(1-\mu)\mu^s]
  \le c_2\int_0^1(1-u)u^s\,du.
\]
The two beta integrals are $(s+1)^{-1}$ and
$[(s+1)(s+2)]^{-1}$.  Dividing \cref{eq:qs} by the lower bound in
\cref{eq:ps} proves \cref{eq:posterior-success-mean}.
\end{proof}

\subsection{The capped sieve}

Fix a target $\eps\in(0,1/8)$ and put $L:=\log n$.  Define
\begin{equation}\label{eq:aj-Sj}
  \begin{aligned}
    a_j&:=\frac{2^j-1}{2^w-1},
       \quad j=0,\ldots,w,\qquad S_0:=0,\\
    S_j&:=\max\left\{S_{j-1}+1,
                \left\lceil\eps^{-a_j}\right\rceil\right\},
       \quad j=1,\ldots,w.
  \end{aligned}
\end{equation}
Thus $a_0=0$, $a_w=1$, and the depths are strictly increasing. For fixed
$w$ and each $j\ge1$, $S_j\asymp_w\eps^{-a_j}$ uniformly over
$\eps\in(0,1/8)$.  Let
\begin{equation}\label{eq:m-def}
  m:=\left\lceil C_0\frac{L}{\eps}\right\rceil
\end{equation}
for a sufficiently large constant $C_0=C_0(w,c_1)$.

Order the arms within each cohort by a fixed label, independent of their
means, and select its first $m$ arms as candidates. A candidate survives
level $j$ if its first $S_j$ \emph{designated screening observations} are
all successes. At age $j-1$, every surviving candidate is scheduled for
\[
  \Delta_j:=S_j-S_{j-1}
\]
additional screening pulls. The policy ignores all nonscreening rewards
when updating this state. In particular, the rule does not stop early
after a failure within a batch.

Reserve $B:=\lfloor n/(2w)\rfloor$ plays for each level and write
$u_j:=\lfloor B/\Delta_j\rfloor$. If more than $u_j$ candidates reach
level $j$, process the $u_j$ smallest labels and permanently discard
the others. We call this an overflow. There is at most one cohort at
each screening age, so the $w$ caps limit total screening to $n/2$ plays.
Assign every remaining play to the smallest-label final survivor of
the age-$w$ cohort. If there is no survivor, use the first arm of
that same cohort as a fallback.

The sieve starts with the first cohort. During the initial $w$ rounds,
only the existing levels are processed, and a fixed alive arm receives
the remaining plays. These finitely many rounds do not affect the
long-run criterion. Ignoring nonscreening rewards ensures that these
initial allocations do not alter the screening experiment.

\begin{algorithm}[b]
\caption{Capped all-success sieve in round $t$}
\label{alg:sieve}
\small
\begin{enumerate}[label=\arabic*.,leftmargin=1.7em,itemsep=5pt]
\item Select the first $m$ arms of $A_t$ as level-$0$ candidates.
\item \textbf{Schedule screening.} For each $j=1,\ldots,\min\{w,t\}$,
let $C_{t,j}$ be the level-$(j-1)$ survivors of $A_{t-j+1}$.
Keep the first $\min\{|C_{t,j}|,u_j\}$ labels and permanently discard
all others. Schedule $\Delta_j$ screening pulls for every kept arm.
Let $Q$ be the total number of screening pulls scheduled across levels.
\item \textbf{Complete the allocation.} If $t>w$, choose the smallest-label
final survivor in $A_{t-w}$, falling back to the first arm of that cohort
if none exists. If $t\le w$, choose a fixed alive arm.
Assign all remaining $n-Q$ plays to this arm.
\item \textbf{Receive feedback.} Reveal all round-$t$ rewards only after
all $n$ plays have been allocated. For each processed cohort, retain
an arm at its next level exactly when all its designated screening
rewards so far equal $1$. Ignore nonscreening rewards.
\end{enumerate}
\end{algorithm}

For the analysis, couple this policy to an ideal sieve with the same
candidates and screening reward streams but no caps. Conditional on
the arm means, these streams consist of independent Bernoulli samples.
Capping removes candidates and their future screening pulls; it never
creates a pull absent from the ideal process.

\begin{lemma}[Equalized endpoint cost]\label{lem:equalized-cost}
For every level $j\in\{1,\ldots,w\}$, the expected endpoint shortfall incurred
by the ideal level-$j$ exploration of one cohort is at most
\begin{equation}\label{eq:stage-cost}
  C\,m\frac{S_j}{(S_{j-1}+1)^2}
  \le C' L\eps^{-\beta_w},
  \qquad
  \beta_w=\frac{2^w}{2^w-1}.
\end{equation}
\end{lemma}

\begin{proof}
In the ideal sieve, a fixed candidate reaches level $j$ exactly when its
first $S_{j-1}$ screening rewards are all successes.  Conditional on its mean, that event has probability
$\mu^{S_{j-1}}$.  Its expected endpoint cost at level $j$ is therefore
\[
  \Delta_j\E[(1-\mu)\mu^{S_{j-1}}]
  \le C\frac{S_j}{(S_{j-1}+1)^2}
\]
by \cref{lem:success-moments}.  Multiply by $m$.

Ignoring ceilings, the exponent is
\[
  1+a_j-2a_{j-1}
  =1+\frac{2^j-1-2(2^{j-1}-1)}{2^w-1}
  =\frac{2^w}{2^w-1}
  =\beta_w.
\]
The recursive rounding in \cref{eq:aj-Sj} changes only a constant depending
on fixed $w$: $S_j\le \lceil\eps^{-a_j}\rceil+j$ and
$S_{j-1}+1\ge\eps^{-a_{j-1}}$.  Since $m\asymp L\eps^{-1}$,
\cref{eq:stage-cost} follows.
\end{proof}

\subsection{Candidate availability and capacity}

The caps make feasibility deterministic. The next lemma shows that the
candidate shortage and cap events have negligible probability at the chosen
sampling scale.

\begin{lemma}[No-survivor and overflow probabilities]\label{lem:upper-feasibility}
There is a constant $C_1=C_1(w,c_1,c_2)$ such that, whenever
\begin{equation}\label{eq:eps-condition}
  \eps\ge C_1\max\left\{\frac{L}{k},
       \left(\frac{L}{n}\right)^{\alpha_w}\right\},
\end{equation}
the following hold for all sufficiently large $n$:
\begin{enumerate}[label=(\alph*),leftmargin=2em]
  \item $m\le k$;
  \item an ideal candidate set has no level-$w$ survivor with probability at
        most $n^{-8}$;
  \item in any fixed global round, the probability that any of the $w$ level
        caps is active is at most $n^{-6}$.
\end{enumerate}
\end{lemma}

\begin{proof}
Part (a) follows from $m\le C L/\eps$ and the first term in
\cref{eq:eps-condition}, after increasing $C_1$.

For part (b), $S_w\le\lceil\eps^{-1}\rceil+w$, and
\cref{eq:ps} gives $p_{S_w}\ge c_w\eps$.  Candidate survival events are
independent across arms, so
\[
  \Pp(\text{no final survivor})
  =(1-p_{S_w})^m
  \le \exp(-mp_{S_w})
  \le n^{-8}
\]
when $C_0$ is sufficiently large.

For part (c), first consider level $1$.  Its demand is deterministic and bounded
by
\begin{equation}\label{eq:stage-one-demand}
  m\Delta_1
  \le C L\eps^{-d_1},
  \qquad d_1:=1+a_1-a_0.
\end{equation}
For a general level $j$, let $Z_{j-1}$ be the number of ideal candidates that
survive $S_{j-1}$ successes.  Then
\[
  Z_{j-1}\sim\Bin(m,p_{S_{j-1}}),
  \qquad
  \lambda_{j-1}:=\E Z_{j-1}
  \le C\frac{m}{S_{j-1}+1}.
\]
Its ideal pull demand is $Q_j=\Delta_jZ_{j-1}$, and
\begin{equation}\label{eq:mean-pull-demand}
  \E Q_j
  \le C L\eps^{-d_j},
  \qquad
  d_j:=1+a_j-a_{j-1}.
\end{equation}
The largest $d_j$ occurs at $j=w$, and
\begin{equation}\label{eq:capacity-margin}
  \alpha_w d_j
  \le \alpha_w d_w
  =\frac{3\cdot2^{w-1}-1}{2^{w+1}-1}
  <1.
\end{equation}
Under \cref{eq:eps-condition}, \cref{eq:mean-pull-demand} is therefore
$o(n)$ uniformly over the fixed set of levels.  In particular,
\cref{eq:stage-one-demand} is at most $B$ for all large $n$.

For $j\ge2$, put $u_j:=\lfloor B/\Delta_j\rfloor$.  The same exponent margin
implies both $u_j\ge 2e\lambda_{j-1}$ and $u_j\ge 12\log n$ for all large $n$.
The standard binomial tail bound gives
\[
  \Pp(Q_j>B)
  =\Pp(Z_{j-1}>u_j)
  \le\left(\frac{e\lambda_{j-1}}{u_j}\right)^{u_j}
  \le 2^{-u_j}
  \le n^{-8}.
\]
A union bound over fixed $w$ gives at most $wn^{-8}\le n^{-6}$.
For each cohort and level, actual survivors form a subset of ideal survivors.
Consequently every actual overflow is contained in the corresponding ideal
overflow event, even when an earlier cap has already discarded candidates.
\end{proof}

\begin{proposition}[Upper bound]\label{prop:upper}
Under \cref{ass:bounded-density}, the capped sieve with any
$\eps\in(0,1/8)$ satisfying \cref{eq:eps-condition} obeys
\begin{equation}\label{eq:upper-eps}
  \overline{\Loss}_n
  \le C_{w,c_1,c_2}\left[
      \eps+\frac{L}{n}\eps^{-\beta_w}+n^{-5}
  \right].
\end{equation}
Consequently, with
\begin{equation}\label{eq:eps-choice}
  \eps=C_1\max\left\{\frac{L}{k},
       \left(\frac{L}{n}\right)^{\alpha_w}\right\},
\end{equation}
we obtain \cref{eq:main-upper}.
\end{proposition}

\begin{proof}
The true loss to the best alive arm is no larger than the endpoint shortfall
$1-\mu$ of the selected arm.  Capping can only remove exploration pulls, so
\cref{lem:equalized-cost}, summed over the fixed $w$ levels and divided by $n$,
bounds exploration loss by
$C_w(L/n)\eps^{-\beta_w}$.

Let $I$ be the smallest-label ideal final survivor of the cohort used for
exploitation.
Conditional on $I=i$, arm $i$ has exactly the prior conditioned on $S_w$
successes: the event that lower-index arms fail to survive is independent of
arm $i$'s mean and rewards.  Hence, without conditioning on any capacity event,
\[
  \E[X_I\1\{I\text{ exists}\}]\le C/S_w\le C_w\eps
\]
by \cref{eq:posterior-success-mean}. Since at most $n$ exploitation plays
are used, this expectation bounds their normalized endpoint cost whenever
the actual and ideal choices agree. The actual exploitation arm equals $I$
unless the ideal cohort has no final survivor or a cap was active during one of
that cohort's $w$ exploration rounds.  A union bound and
\cref{lem:upper-feasibility} put this mismatch probability at $O_w(n^{-6})$;
on the mismatch event the per-pull shortfall is at most one.  This proves
\cref{eq:upper-eps}.

Finally, $\alpha_w(1+\beta_w)=1$.  Hence
$\eps\ge(L/n)^{\alpha_w}$ implies
\[
  \frac{L}{n}\eps^{-\beta_w}\le\eps.
\]
Substituting \cref{eq:eps-choice} gives \cref{eq:main-upper}.  If the chosen
$\eps$ is not below $1/8$, the desired upper bound is already a constant and the
policy that always plays the first arm of the fresh cohort suffices after
enlarging $C_{w,c_1,c_2}$.
\end{proof}

\section{A batched posterior-discovery bound}\label{sec:discovery}

We now establish the information-theoretic core of the lower bound.  Throughout
this section, $X=1-\mu\sim\Unif[0,1]$.  A reward success has probability $1-X$
and a failure has probability $X$.  After $F$ failures and $S$ successes,
\[
  X\mid\F\sim\BetaDist(F+1,S+1),
  \qquad
  g:=\E[X\mid\F]=\frac{F+1}{F+S+2}.
\]

\subsection{Posterior small balls and hitting}

\begin{lemma}[Beta posterior small ball]\label{lem:beta-small-ball}
Let $X\sim\BetaDist(a,b)$ with $a,b\ge1$ and mean $g=a/(a+b)$.  For every
$y\in(0,1)$,
\begin{equation}\label{eq:beta-small-ball}
  \Pp(X\le y)
  \le e\min\left\{1,\frac{y}{g}\right\}.
\end{equation}
\end{lemma}

\begin{proof}
The bound is trivial when $y\ge g$.  Suppose $y<g$.  Write
$X=U/(U+V)$ for independent $U\sim\operatorname{Gamma}(a,1)$ and
$V\sim\operatorname{Gamma}(b,1)$.  For
$0<\lambda<1/y$, Chernoff's inequality gives
\[
  \Pp(X\le y)
  =\Pp((1-y)U-yV\le0)
  \le (1+\lambda(1-y))^{-a}(1-\lambda y)^{-b}.
\]
At
\[
  \lambda_*:=\frac{g-y}{y(1-y)},
\]
the right-hand side equals
\begin{equation}\label{eq:beta-chernoff}
  \left(\frac{y}{g}\right)^a
  \left(\frac{1-y}{1-g}\right)^b.
\end{equation}
Set $r=y/g$.  Since
\[
  b\log\frac{1-y}{1-g}
  \le \frac{b(g-y)}{1-g}
  =a(1-r),
\]
\cref{eq:beta-chernoff} is at most
$[r e^{1-r}]^a\le r e^{1-r}\le er$.
\end{proof}

\begin{lemma}[Posterior hitting probability]\label{lem:posterior-hitting}
Start from a $\BetaDist(a,b)$ posterior with $a,b\ge1$ and mean $g$.  Under any nonanticipating future
sampling rule, observed at any finite collection of future batch endpoints, and
for $x\in(0,1/4)$,
\begin{equation}\label{eq:hitting}
  \Pp\!\left(\min_r g_r\le x\mid\F\right)
  \le C\min\left\{1,\frac{x}{g}\right\}.
\end{equation}
The rule may depend on arbitrary auxiliary observations whose conditional law,
given the arm's history, contains no additional information about its latent
mean.
\end{lemma}

\begin{proof}
Let $\tau$ be the first permitted update time with $g_\tau\le x$.  On
$\{\tau<\infty\}$, Markov's inequality under the posterior at time $\tau$
gives
\[
  \Pp(X\le2x\mid\F_\tau)\ge\frac12.
\]
Because $\{\tau<\infty\}\in\F_\tau$, the tower property yields
\[
  \Pp(\tau<\infty\mid\F)
  \le2\Pp(X\le2x\mid\F).
\]
Apply \cref{lem:beta-small-ball} and absorb constants.  Predictable adaptive
sampling and conditionally uninformative auxiliary variables leave the beta
posterior likelihood unchanged, so the same argument applies in the enlarged
filtration.
\end{proof}

\subsection{The discovery experiment}

Consider a finite set of target arms with independent uniform priors.
They are sampled in $j$ batches, with arbitrary nonnegative integer
counts and finite expected total cost. Write $\Hh_{r-1}^{-}$ for the
history just before choosing the batch-$r$ counts, and
$\Hh_{r-1}^{+}$ for the history including the independent randomization
and chosen counts $q_{i,r}$. Thus $q_{i,r}$ is
$\Hh_{r-1}^{+}$-measurable, and the target posteriors are unchanged
between these two sigma-fields. The batch rewards are revealed together
after all counts have been chosen.

The history may contain observations from independent outside arms.
The relevant requirement is posterior factorization under the full
adaptive transcript: outside observations and predictable choices
introduce no additional likelihood factor in a target mean.
\Cref{app:external} proves this property for the streaming experiment.

Let $g_{i,r}$ be arm $i$'s posterior mean after batch $r$, with
$g_{i,0}=1/2$, and define
\[
  M_{i,j}:=\min_{0\le r\le j}g_{i,r},
  \qquad
  H_j(x):=\E\sum_i\1\{M_{i,j}\le x\}.
\]
The predictable Bayesian endpoint cost through batch $j$ is
\begin{equation}\label{eq:Bj}
  B_j:=\E\sum_{r=1}^{j}\sum_i q_{i,r}g_{i,r-1}.
\end{equation}
By the tower property, this is exactly the expected sum of true endpoint
shortfalls over all target-arm samples.

\begin{lemma}[One-step discovery recursion]\label{lem:discovery-recursion}
For $0<x<\Gamma<1/4$,
\begin{align}
  H_j(x)
  \le{}&H_{j-1}(x)
  +Cx\frac{H_{j-1}(\Gamma)}{\Gamma}
  +Cx\int_x^\Gamma\frac{H_{j-1}(s)}{s^2}\,ds
  +CB_j\frac{x^2}{\Gamma^2}.                 \label{eq:discovery-recursion}
\end{align}
\end{lemma}

\begin{proof}
We split arms by $M_{i,j-1}$.

Arms with $M_{i,j-1}\le x$ contribute $H_{j-1}(x)$.  For an arm with
$x<M_{i,j-1}\le\Gamma$, its current posterior mean satisfies
$g_{i,j-1}\ge M_{i,j-1}$.  Conditional on $\Hh_{j-1}^{+}$,
\cref{lem:posterior-hitting} bounds its chance of crossing $x$ in batch $j$ by
$Cx/M_{i,j-1}$.  Therefore the contribution of these arms is at most
\[
  Cx\E\sum_i
  \frac{\1\{x<M_{i,j-1}\le\Gamma\}}{M_{i,j-1}}.
\]
Stieltjes integration by parts for the counting measure gives
\[
  \E\sum_i
  \frac{\1\{x<M_{i,j-1}\le\Gamma\}}{M_{i,j-1}}
  \le \frac{H_{j-1}(\Gamma)}{\Gamma}
  +\int_x^\Gamma\frac{H_{j-1}(s)}{s^2}\,ds,
\]
where the omitted boundary term is nonpositive.  This yields the first three
terms in \cref{eq:discovery-recursion}.

It remains to handle arms with $M_{i,j-1}>\Gamma$.  Let $N_{i,j}$ be the total
number of samples of arm $i$ through batch $j$.  A beta posterior after
$N_{i,j}$ observations has mean at least $1/(N_{i,j}+2)$.  Thus
$g_{i,j}\le x$ and $x<1/4$ imply
\[
  N_{i,j}\ge \frac{1}{2x}.
\]
Define the $\Hh_{j-1}^{+}$-measurable event
\[
  E_i:=\left\{M_{i,j-1}>\Gamma,
              \ N_{i,j}\ge\frac{1}{2x}\right\}.
\]
On $E_i$, every sample of arm $i$ was allocated at a batch start whose posterior
mean exceeded $\Gamma$.  Hence its pathwise predictable cost
\[
  C_i:=\sum_{r=1}^{j}q_{i,r}g_{i,r-1}
\]
satisfies $C_i\ge\Gamma/(2x)$, and therefore
\begin{equation}\label{eq:heavy-count}
  \sum_i\1_{E_i}
  \le \frac{2x}{\Gamma}\sum_i C_i.
\end{equation}
Conditional on $\Hh_{j-1}^{+}$ and $E_i$, the current posterior mean still
exceeds $\Gamma$, so \cref{lem:posterior-hitting} gives
\[
  \Pp(g_{i,j}\le x\mid\Hh_{j-1}^{+})
  \le C\frac{x}{\Gamma}.
\]
Multiplying this conditional bound by the measurable indicator $\1_{E_i}$,
summing, taking expectations, and using \cref{eq:heavy-count} yields
\[
  \E\sum_i\1_{E_i}\1\{g_{i,j}\le x\}
  \le C\frac{x}{\Gamma}\E\sum_i\1_{E_i}
  \le CB_j\frac{x^2}{\Gamma^2}.
\]
This is the final term.
\end{proof}

\begin{theorem}[Batched posterior-discovery inequality]\label{thm:batched-discovery}
For every fixed integer $j\ge1$, there are constants $C_j<\infty$ and
$x_j>0$ such that every $j$-batch experiment described above satisfies
\begin{equation}\label{eq:discovery-bound}
  H_j(x)\le C_j B_j x^{\beta_j},
  \qquad
  \beta_j:=\frac{2^j}{2^j-1},
  \qquad 0<x\le x_j.
\end{equation}
Consequently,
\begin{equation}\label{eq:discovery-prob}
  \Pp\!\left(\exists i:\ M_{i,j}\le x\right)
  \le C_j B_jx^{\beta_j}.
\end{equation}
\end{theorem}

\begin{proof}
For $j=1$, choose a fixed $\Gamma\in(0,1/4)$ in
\cref{lem:discovery-recursion}.  Since every fresh arm has prior mean $1/2$,
$H_0(\Gamma)=0$, and
\[
  H_1(x)\le C B_1x^2.
\]
Thus $\beta_1=2$.

Suppose for some $\beta>1$ that
$H_{j-1}(s)\le A B_{j-1}s^\beta$ throughout a sufficiently small interval.
Since $B_{j-1}\le B_j$, \cref{lem:discovery-recursion} gives
\[
  H_j(x)
  \le C A B_j x\Gamma^{\beta-1}
       +C B_j\frac{x^2}{\Gamma^2}.
\]
The integral contributes a factor $1/(\beta-1)$, which is absorbed
into the constant at this fixed induction step. The term $H_{j-1}(x)$
is also absorbed because
$x^\beta\le x\Gamma^{\beta-1}$ for $x<\Gamma$.  Choose
\[
  \Gamma=x^{1/(\beta+1)}.
\]
For sufficiently small $x$, this lies in $(x,1/4)$, and
\[
  x\Gamma^{\beta-1}
  =\frac{x^2}{\Gamma^2}
  =x^{2\beta/(\beta+1)}.
\]
Thus
\[
  \beta_j=\frac{2\beta_{j-1}}{\beta_{j-1}+1},
  \qquad \beta_1=2.
\]
Equivalently,
$1/\beta_j=(1+1/\beta_{j-1})/2$, so
$1/\beta_j=1-2^{-j}$ and
$\beta_j=2^j/(2^j-1)$.  Constants and admissible radii can be chosen
inductively because $j$ is fixed.  Finally,
\cref{eq:discovery-prob} follows from
$\Pp(\sum_i\1\{M_{i,j}\le x\}\ge1)\le H_j(x)$.
\end{proof}

\section{Embedding discovery into streaming cohorts}\label{sec:embedding}

We now return to SLHVB under the uniform prior.  Let
\[
  g_{a,t}:=\E[X_a\mid\F_{t-1}],
  \qquad
  G_t:=\min_{a\in\A_t}g_{a,t}.
\]
Define the per-round endpoint shortfall
\begin{equation}\label{eq:Rt}
  R_t:=\frac1n\E\sum_{a\in\A_t}\pi_t(a)X_a
      =\frac1n\E\sum_{a\in\A_t}\pi_t(a)g_{a,t}.
\end{equation}
The second equality is the tower property.  Since a weighted average is no
smaller than its minimum,
\begin{equation}\label{eq:G-vs-R}
  \E G_t\le R_t.
\end{equation}

\begin{lemma}[Fixed-round streaming discovery]\label{lem:streaming-discovery}
For each $t>w$, let
\begin{equation}\label{eq:Bt}
  B_t:=\E\sum_{s=t-w}^{t-1}\sum_{a\in\A_t}
           \pi_s(a)g_{a,s},
\end{equation}
where $\pi_s(a)=0$ before arm $a$ is born.  For all sufficiently small $x$,
\begin{equation}\label{eq:fixed-round-discovery}
  \Pp(G_t\le x)
  \le C_w B_t x^{\beta_w},
\end{equation}
and
\begin{equation}\label{eq:Bt-upper}
  B_t\le n\sum_{s=t-w}^{t-1}R_s.
\end{equation}
\end{lemma}

\begin{proof}
Fix $t$ and regard the target arms as the union of cohorts
$A_{t-w},\ldots,A_t$.  Their histories before round $t$ can be exposed in the
$w$ global batches $t-w,\ldots,t-1$, padding an arm's allocation by zero before
its birth.  All target arms are fresh at the beginning of this window; later cohorts
can equivalently be drawn in advance and left unobserved until arrival.
Before-window information is independent of their means. Rewards from
outside arms may affect allocations, but posterior factorization under
the full transcript, proved in \cref{app:external}, makes them admissible
auxiliary observations in \cref{thm:batched-discovery}.

If $G_t\le x$, at least one target arm has current posterior mean at most $x$,
and therefore its minimum posterior mean over the $w$ batch endpoints is at
most $x$.  Applying \cref{eq:discovery-prob} gives
\cref{eq:fixed-round-discovery}.  Finally, the target-arm sum in each past round
is bounded by the sum over all arms, whose expectation is $nR_s$ by
\cref{eq:Rt}.  This proves \cref{eq:Bt-upper}.
\end{proof}

\begin{proposition}[Endpoint-shortfall lower bound]\label{prop:endpoint-lower}
Under the uniform prior, every feasible policy satisfies
\begin{equation}\label{eq:endpoint-lower}
  \liminf_{T\to\infty}\frac1T\sum_{t=1}^{T}R_t
  \ge c_w n^{-\alpha_w}.
\end{equation}
\end{proposition}

\begin{proof}
Put
\[
  \overline R_T:=\frac1T\sum_{t=1}^{T}R_t
\]
and let $U$ be uniform on $\{w+1,\ldots,T\}$, independently of all other
randomness.  For $T\ge2w$, \cref{eq:G-vs-R} gives
\[
  \E G_U
  \le \frac{T}{T-w}\overline R_T
  \le2\overline R_T.
\]
Markov's inequality implies
\begin{equation}\label{eq:markov-good-arm}
  \Pp(G_U\le4\overline R_T)\ge\frac12.
\end{equation}
If $4\overline R_T$ exceeds the small constant $x_w$ from
\cref{thm:batched-discovery}, then
$\overline R_T\ge x_w/4$, which is stronger than the desired claim for all
large $n$.  Otherwise, average \cref{eq:fixed-round-discovery} over $t$ and use
\cref{eq:Bt-upper}:
\begin{align*}
  \frac12
  &\le \frac1{T-w}\sum_{t=w+1}^{T}
        \Pp(G_t\le4\overline R_T) \\
  &\le C_w\overline R_T^{\beta_w}
        \frac1{T-w}\sum_{t=w+1}^{T}B_t \\
  &\le C_w\overline R_T^{\beta_w}
        \frac{n}{T-w}
        \sum_{t=w+1}^{T}\sum_{s=t-w}^{t-1}R_s \\
  &\le C_w n\overline R_T^{\beta_w+1}.
\end{align*}
The last step uses the fact that each $R_s$ occurs in at most $w$ windows and
$T/(T-w)\le2$, with fixed factors absorbed into $C_w$.  Hence
\[
  \overline R_T
  \ge c_w n^{-1/(1+\beta_w)}
  =c_w n^{-\alpha_w}.
\]
This holds uniformly for all sufficiently large $T$, proving the liminf bound.
\end{proof}

\section{The arrival floor and the loss lower bound}\label{sec:arrival-floor}

The endpoint quantity $R_t$ differs from actual loss by the shortfall of the
best alive arm.  To complete the phase diagram, we also need a $1/k$ lower bound.
We give a direct proof under the uniform prior.

\begin{lemma}[An explicit arrival floor]\label{lem:arrival-floor}
For $D=\Unif[0,1]$, every admissible policy, every $k\ge1$, and
every $t>w$ satisfy
\begin{equation}\label{eq:arrival-floor}
  L_t:=\frac1n\E\sum_{a\in\A_t}
       \pi_t(a)(\mu_t^*-\mu_a)
  \ge \frac{\eta_w}{k},
  \qquad
  \eta_w:=\frac{(3/4)^w-(1/2)^w}{64}>0.
\end{equation}
\end{lemma}

\begin{proof}
Let $Y$ be the largest mean among the $wk$ old arms and $Z$ the
largest mean among the $k$ fresh arms:
\[
  Y:=\max_{a\in A_{t-w}\cup\cdots\cup A_{t-1}}\mu_a,
  \qquad Z:=\max_{a\in A_t}\mu_a.
\]
Write $N_t:=\sum_{a\in A_t}\pi_t(a)$ for the number of plays of fresh arms.
Since no old arm exceeds $Y$ and no fresh arm exceeds $Z$, the
round's regret is at least
\begin{equation}\label{eq:pathwise-arrival}
  N_t(Y-Z)_+ +(n-N_t)(Z-Y)_+.
\end{equation}
Condition on all old means, the past history, and the randomization
used for the current allocation. Then $Y$ and $N_t$ are fixed, and
the fresh means remain independent uniform variables.
For $q=N_t/n$, the conditional expected loss is at least
$qA_k(Y)+(1-q)B_k(Y)\ge\min\{A_k(Y),B_k(Y)\}$, where
\begin{align}
  A_k(y)&:=\E[(y-Z)_+]=\frac{y^{k+1}}{k+1}, \label{eq:Ak}\\
  B_k(y)&:=\E[(Z-y)_+]
       =\int_y^1(1-u^k)\,du
       =(1-y)-\frac{1-y^{k+1}}{k+1}. \label{eq:Bk}
\end{align}
These formulas follow directly from $\Pp(Z\le z)=z^k$.

Consider the event
\[
  \mathcal W:=\{1/2\le Y^k\le3/4\}.
\]
Because $\Pp(Y\le y)=y^{wk}$,
\[
  \Pp(\mathcal W)=(3/4)^w-(1/2)^w.
\]
On this event, $Y^k\ge1/2$ and $Y\ge1/2$, so
$A_k(Y)\ge1/[4(k+1)]\ge1/(8k)$. Also
$Y\le(3/4)^{1/k}$, which gives
\begin{align*}
  B_k(Y)
  &\ge\int_{(3/4)^{1/k}}^{(7/8)^{1/k}}(1-u^k)\,du\\
  &\ge\frac18\left[(7/8)^{1/k}-(3/4)^{1/k}\right]
   \ge\frac{1}{64k}.
\end{align*}
The last inequality uses $(v^{1/k})'\ge1/k$ for $0<v<1$.
Thus the smaller of $A_k(Y)$ and $B_k(Y)$ is at least $1/(64k)$
on $\mathcal W$.
Taking expectations proves the claim for every $k\ge1$.
\end{proof}

For the uniform prior, after warm-up there are $(w+1)k$ alive arms, so the
minimum endpoint shortfall satisfies
\begin{equation}\label{eq:best-endpoint}
  \E X_t^*=\frac{1}{(w+1)k+1}\le\frac{1}{k},
  \qquad
  X_t^*:=1-\mu_t^*.
\end{equation}
Because exactly $n$ plays are used,
\begin{equation}\label{eq:R-L-relation}
  R_t=L_t+\E X_t^*,
\end{equation}
where $L_t$ is the expected per-pull loss in round $t$.

\begin{proof}[Proof of the lower bound in \cref{thm:main}]
Put $L:=\liminf_{T\to\infty}T^{-1}\sum_{t=1}^{T}L_t$.
The endpoint lower bound and \cref{eq:best-endpoint,eq:R-L-relation}
give
\begin{equation}\label{eq:lower-combination-one}
  L+\frac1k\ge b_w n^{-\alpha_w}
\end{equation}
for some $b_w>0$. Meanwhile, \cref{lem:arrival-floor} gives
\begin{equation}\label{eq:lower-combination-two}
  L\ge\frac{\eta_w}{k}.
\end{equation}
Hence $1/k\le L/\eta_w$, and substitution in
\cref{eq:lower-combination-one} yields
\[
  L\ge\frac{\eta_w b_w}{1+\eta_w}\,n^{-\alpha_w}.
\]
Combining this with \cref{eq:lower-combination-two} proves
$L\ge c_w\max\{k^{-1},n^{-\alpha_w}\}$.
Together with \cref{prop:upper}, this completes the proof of
\cref{thm:main}.
\end{proof}

\section{Consequences and scope}\label{sec:discussion}

\paragraph{The arrival threshold.}
\Cref{tab:regimes} separates the part of the parameter space that was
already understood from the region settled by the present bounds.
The range in which loss is determined by the arrival rate expands
from $\rho\le\theta_w$ to $\rho\le\alpha_w$. Above the latter
threshold, additional arrivals do not improve the polynomial rate:
the number of feedback opportunities is the binding constraint.
The upper and lower bounds match at the exponent level throughout;
the logarithmic factors in \cref{eq:main-upper} remain unresolved.

\paragraph{The cost of finite feedback.}
For every fixed $w\ge1$,
\[
  \theta_w<\alpha_w<\frac12,
  \qquad
  \alpha_w-\theta_w
  =\frac{2^{w+1}-(w+2)}
         {2(w+1)(2^{w+1}-1)}.
\]
The previous upper bound is therefore polynomially improvable whenever
its lifetime term is active. Conversely, the uniform-prior lower bound
shows that no fixed lifetime attains exponent $1/2$ in the high-arrival
regime. By \cref{eq:exponent-gap}, each additional feedback opportunity
approximately halves the remaining exponent gap. This conclusion
concerns the sequence of fixed-$w$ exponents; it does not assert a
bound uniform over $w=w(n)$.

\paragraph{What Bernoulli observations contribute.}
Near $\mu=1$, the failure probability and the variance
$\mu(1-\mu)$ both vanish. An all-success event therefore selects
rare good arms and controls their posterior shortfall simultaneously.
Variance-independent confidence bounds do not exploit this feature.
The present upper bound uses it through the two moment laws
$p_s\asymp(s+1)^{-1}$ and $q_s\lesssim(s+1)^{-2}$.
The lower bound shows that no other admissible policy improves the
polynomial exponent under the uniform prior.

\paragraph{Implementation.}
The policy maintains candidate identities, ages, and whether each
completed screening batch contained a failure. Its stage depths and
caps are computed in advance, and it needs neither dynamic programming
nor a posterior index. Every scheduled pull is executed; an arm is
discarded only after the end of the first screening batch containing
a failure, or when a cap removes it. These features establish a
simple implementable policy, although the present paper makes no
empirical performance claim.

\paragraph{Open directions.}
Three questions remain beyond the theorem. First, the logarithmic
gap may be reduced by a more economical candidate-selection rule.
Second, a matching pointwise lower bound for every bounded-density
prior would require extending the posterior argument beyond the
uniform prior used here. Third, other endpoint tails and observation
families can change both the screening moments and posterior
small-ball behavior. In particular, the present theorem does not
determine the optimal finite-lifetime rate for homoscedastic
sub-Gaussian observations.

\section{Conclusion}

The optimal loss exponent for Bernoulli SLHVB with fixed lifetime
parameter $w$ is
$\alpha_w=(2^w-1)/(2^{w+1}-1)$.
A capped all-success sieve achieves this exponent for every
bounded-density prior, and a posterior-discovery inequality proves
its necessity under the uniform prior. The resulting minimax rate
separates the unavoidable cost of new arrivals from the cost of
limited feedback and closes the previous polynomial gap in this
Bernoulli model.

\appendix

\section{Depth identities and capacity margins}\label{app:algebra}

The sampling depths in \cref{eq:aj-Sj} satisfy
\begin{align*}
  1+a_j-2a_{j-1}&=\frac{2^w}{2^w-1}=\beta_w,\\
  d_j:=1+a_j-a_{j-1}&=1+\frac{2^{j-1}}{2^w-1},\\
  1-\alpha_w\max_{1\le j\le w}d_j
    &=\frac{2^{w-1}}{2^{w+1}-1}>0,\\
  \alpha_w(1+\beta_w)&=1.
\end{align*}
The third identity is the capacity margin used in
\cref{lem:upper-feasibility}. With $L=\log n$ and
$\eps\ge C(L/n)^{\alpha_w}$,
\[
  \frac{\E Q_j}{n}
  \le C\frac{L}{n}\eps^{-d_j}
  \le C n^{\alpha_wd_j-1}L^{1-\alpha_wd_j}
  =o(1)
\]
uniformly over the fixed set of levels and all admissible $\eps$.
For sufficiently large $n$, $B\ge n/(4w)$ and
$\Delta_j\le C_w\eps^{-a_j}$. It follows that
\[
  u_j=\left\lfloor\frac{B}{\Delta_j}\right\rfloor
  \ge c_w n\eps^{a_j}
  \ge c_w n^{1-\alpha_wa_j}L^{\alpha_wa_j}.
\]
Because $a_j\le1$ and $\alpha_w<1/2$, this quantity dominates
$\log n$. In addition,
$u_j/\lambda_{j-1}$ is, up to fixed constants, the reciprocal of
$\E Q_j/B$ and tends to infinity. Therefore both
$u_j\ge2e\lambda_{j-1}$ and $u_j\ge12\log n$ hold for all large $n$.
The threshold on $n$ depends on $w,c_1,c_2$, but not on $k$.

\section{Posterior factorization under adaptive sampling}\label{app:external}

We justify the treatment of independent outside observations in
\cref{thm:batched-discovery,lem:streaming-discovery}.
Let $X_1,\ldots,X_M$ be the target shortfalls. They have independent
uniform priors. All outside latent variables and the policy's random
seed are independent of them.

Fix a finite realized transcript $h$, including allocations and all
observed rewards. Write $F_i(h)$ and $S_i(h)$ for the failures and
successes observed from target arm $i$. The likelihood of this
transcript, after integrating out outside latent variables, has the form
\begin{equation}\label{eq:likelihood-factorization}
  \mathcal L_h(x_1,\ldots,x_M)
  =C(h)\prod_{i=1}^{M}x_i^{F_i(h)}(1-x_i)^{S_i(h)},
\end{equation}
where $C(h)$ does not depend on the target shortfalls.
To see this, factor the transcript into its successive decision and
observation kernels. Conditional on the preceding transcript, a
nonanticipating decision kernel contains no unknown target mean.
Target observations contribute the displayed Bernoulli likelihoods.
Outside observations contribute factors involving only outside
latent variables and the fixed transcript. Integrating those
variables preserves the stated form. The same argument applies
after an allocation is chosen but before its rewards are revealed.

Bayes' rule now gives independent conditional posteriors
\[
  X_i\mid h\sim\BetaDist(F_i(h)+1,S_i(h)+1).
\]
In particular, even when outside feedback influences future target
allocations, it does not change a target posterior beyond the
likelihood of that arm's own observed rewards. This is the precise
conditional statement used in the discovery recursion.

\section{Long-run criteria and the order of expectation}\label{app:long-run}

The limsup of an expected average and the expectation of a conditional
limsup need not agree for an arbitrary nonstationary policy.
For completeness, the same exponent holds under either convention.
Let $\boldsymbol\mu$ denote the full sequence of latent cohort means and set
\[
  \ell_t(\boldsymbol\mu)
  :=\frac1n\E\!\left[
       \sum_{a\in\A_t}\pi_t(a)(\mu_t^*-\mu_a)
       \,\middle|\,\boldsymbol\mu\right],
  \qquad
  Z_T(\boldsymbol\mu):=\frac1T\sum_{t=1}^{T}\ell_t(\boldsymbol\mu).
\]
Then $0\le Z_T\le1$ and
$\E_D Z_T=\Loss_{n,T}(\pi;D)$.
For the alternative criterion
\[
  \Loss_n^{\mathrm{path}}(\pi;D)
  :=\E_D\!\left[\limsup_{T\to\infty}Z_T(\boldsymbol\mu)\right],
\]
reverse Fatou's inequality yields
\[
  \Loss_n^{\mathrm{path}}(\pi;D)
  \ge \limsup_{T\to\infty}\E_D Z_T
  \ge \underline{\Loss}_n(\pi;D).
\]
Thus the uniform-prior lower bound transfers to this convention for
every admissible policy.

For the upper bound, use the deterministic tie-breaking and the
age-$w$ fallback specified in \cref{alg:sieve}.
After initialization, each cohort's screening state depends only
on its own screening rewards. Therefore
$\ell_t(\boldsymbol\mu)$ is the same bounded measurable function of
the $w+1$ cohort-mean vectors in $A_{t-w},\ldots,A_t$ at every
$t>w$. These vectors are independent and identically distributed
across cohorts. For each residue class modulo $w+1$, the corresponding
values of $\ell_t$ are consequently independent and identically
distributed. The strong law applied to these finitely many
subsequences gives
\[
  Z_T\longrightarrow \E_D\ell_{w+1}
  \quad\text{almost surely}.
\]
Bounded convergence then identifies
$\Loss_n^{\mathrm{path}}=\overline{\Loss}_n=\E_D\ell_{w+1}$
for this policy. The policy that always plays the first fresh arm, used in the trivial
parameter range, has the same finite-window property. Hence
\cref{thm:main,cor:rho} remain valid with
$\Loss_n^{\mathrm{path}}$ in place of the expected-average criterion.

\begin{thebibliography}{13}

\bibitem[Agarwal et~al.(2017)]{agarwal2017limited}
Arpit Agarwal, Shivani Agarwal, Sepehr Assadi, and Sanjeev Khanna.
\newblock Learning with limited rounds of adaptivity: Coin tossing,
multi-armed bandits, and ranking from pairwise comparisons.
\newblock In \emph{Proceedings of the 2017 Conference on Learning Theory},
volume 65 of \emph{Proceedings of Machine Learning Research},
pages 39--75, 2017.
\newblock \url{https://proceedings.mlr.press/v65/agarwal17c.html}.

\bibitem[Bayati et~al.(2020)]{bayati2020greedy}
Mohsen Bayati, Nima Hamidi, Ramesh Johari, and Khashayar Khosravi.
\newblock The unreasonable effectiveness of greedy algorithms in
multi-armed bandit with many arms.
\newblock In \emph{Advances in Neural Information Processing Systems},
volume 33, 2020.

\bibitem[Berry et~al.(1997)]{berry1997infinite}
Donald A. Berry, Robert W. Chen, Alan Zame, David C. Heath, and Larry A. Shepp.
\newblock Bandit problems with infinitely many arms.
\newblock \emph{The Annals of Statistics}, 25(5):2103--2116, 1997.
\newblock \href{https://doi.org/10.1214/aos/1069362389}
  {doi:10.1214/aos/1069362389}.

\bibitem[Bonald and Prouti\`ere(2013)]{bonald2013two}
Thomas Bonald and Alexandre Prouti\`ere.
\newblock Two-target algorithms for infinite-armed bandits with Bernoulli rewards.
\newblock In \emph{Advances in Neural Information Processing Systems},
volume 26, 2013.

\bibitem[Chakrabarti et~al.(2008)]{chakrabarti2008mortal}
Deepayan Chakrabarti, Ravi Kumar, Filip Radlinski, and Eli Upfal.
\newblock Mortal multi-armed bandits.
\newblock In \emph{Advances in Neural Information Processing Systems},
volume 21, pages 273--280, 2008.

\bibitem[Farias and Madan(2011)]{farias2011irrevocable}
Vivek F. Farias and Ritesh Madan.
\newblock The irrevocable multiarmed bandit problem.
\newblock \emph{Operations Research}, 59(2):383--399, 2011.
\newblock \href{https://doi.org/10.1287/opre.1100.0891}
  {doi:10.1287/opre.1100.0891}.

\bibitem[Gao et~al.(2019)]{gao2019batched}
Zijun Gao, Yanjun Han, Zhimei Ren, and Zhengqing Zhou.
\newblock Batched multi-armed bandits problem.
\newblock In \emph{Advances in Neural Information Processing Systems},
volume 32, 2019.

\bibitem[Jia et~al.(2023)]{jia2023short}
Su Jia, Nishant Oli, Ian Anderson, Paul Duff, Andrew A. Li, and R. Ravi.
\newblock Short-lived high-volume bandits.
\newblock In \emph{Proceedings of the 40th International Conference on
Machine Learning}, volume 202 of \emph{Proceedings of Machine Learning
Research}, pages 14902--14929, 2023.
\newblock \url{https://proceedings.mlr.press/v202/jia23b.html}.

\bibitem[Jia et~al.(2026)]{jia2026short}
Su Jia, Andrew Li, R. Ravi, Nishant Oli, Paul Duff, and Ian Anderson.
\newblock Short-lived high-volume bandits.
\newblock \emph{Operations Research}, 74(4):1945--1964, 2026.
\newblock \href{https://doi.org/10.1287/opre.2023.0557}
  {doi:10.1287/opre.2023.0557}.

\bibitem[Levine et~al.(2017)]{levine2017rotting}
Nir Levine, Koby Crammer, and Shie Mannor.
\newblock Rotting bandits.
\newblock In \emph{Advances in Neural Information Processing Systems},
volume 30, 2017.

\bibitem[Perchet et~al.(2016)]{perchet2016batched}
Vianney Perchet, Philippe Rigollet, Sylvain Chassang, and Erik Snowberg.
\newblock Batched bandit problems.
\newblock \emph{The Annals of Statistics}, 44(2):660--681, 2016.
\newblock \href{https://doi.org/10.1214/15-AOS1381}
  {doi:10.1214/15-AOS1381}.

\bibitem[Pilliat(2026)]{pilliat2026adaptive}
Emmanuel Pilliat.
\newblock Adaptive algorithms for infinitely many-armed bandits:
A unified framework.
\newblock \emph{Journal of Machine Learning Research}, 27(192):1--25, 2026.
\newblock \url{https://jmlr.org/papers/v27/25-2669.html}.

\bibitem[Wang et~al.(2008)]{wang2008infinite}
Yizao Wang, Jean-Yves Audibert, and R\'emi Munos.
\newblock Algorithms for infinitely many-armed bandits.
\newblock In \emph{Advances in Neural Information Processing Systems},
volume 21, 2008.

\end{thebibliography}

\end{document}
